\documentclass[11pt,a4paper]{article}
\usepackage[T1]{fontenc}
\usepackage{amsmath,amssymb,geometry}
\geometry{margin=27mm}
\title{Section 3: the complementary five-cycle from one proper profile}
\date{16 October 2026}
\begin{document}
\maketitle
\noindent\textbf{Manuscript policy.}
This is a source-local validation supplement only.
Do not modify the author manuscript.

\section*{An alternative proof for the nondegenerate case of Claim 3.19}
Let \(H\) be a vertex-transitive graph on ten vertices,
\(e:C_5\hookrightarrow H\) an induced EPPA embedding,
and \(B=V(H)\setminus e(V(C_5))\).
Suppose some \(b\in B\) has a nonempty proper
selected neighbourhood \(P\subset V(C_5)\).

By the five-cycle free-orbit lemma, the five selected
neighbourhoods of the vertices in \(B\) are
\(\{P+d:d\in\mathbb Z/5\mathbb Z\}\), each exactly once.
Thus the profile map is injective, and we can canonically
index the actual vertices of \(B\) by \(d\in\mathbb Z/5\mathbb Z\)
via \(P_d=P+d\).

For any rotation \(t\) of the selected five-cycle,
choose an EPPA lift of that single automorphism.
The lift preserves \(B\) and sends selected-neighbourhood
profiles \(P_d\) to \(P_d+t=P_{d+t}\).
Injectivity forces this lift to act on \(B\) by
\(d\mapsto d+t\). No coherent family of extensions is needed.
Consequently the induced graph \(H[B]\), in these canonical
coordinates, is invariant under the cyclic translations.

There are exactly two orbits of the rotation group on
unordered pairs of distinct labels, consecutive pairs and
diagonals. Therefore \(H[B]\) is empty, the standard
\(C_5\), its complementary \(C_5\), or complete.
Vertex-transitivity of \(H\) implies equal numbers of edges
in the two equal halves:
\[
  |E(H[B])|=|E(C_5)|=5.
\]
Thus the empty and complete possibilities are excluded;
\(H[B]\) is \(C_5\) or \(\overline{C_5}\) in the
specified coordinates. Since multiplication by two
identifies \(C_5\) with its complement,
\[
       H[B]\cong C_5.
\]

\paragraph{Scope.}
This argument does not use the imprimitivity of \(H\)
or a prior assertion that \(H[B]\) is vertex-transitive.
It proves the complementary five-cycle only when
at least one outside profile is nonempty and proper.
If all outside profiles are empty or full,
a separate graph-theoretic argument remains necessary.

\section*{Warning: imprimitivity is still essential for the full theorem}
The larger classification must \emph{not} simply drop the
imprimitivity hypothesis on \(H\).
For example, the two-layer graph with bottom \(C_5\),
top \(\overline{C_5}\), and a perfect matching between
equal labels is the Petersen graph. It is transitive,
primitive, and extends all automorphisms of the designated
\(C_5\). The base \(C_5\) is homogeneous, so it is
also an EPPA witness for that copy. This additional
primitive host is not among the imprimitive alternatives
of Lemma~3.18. This observation is a \emph{boundary
check}, not a counterexample to the manuscript lemma.

\section*{Formalisation target}
The new module
\texttt{EPPAIII/BlockSystems/CycleFiveProperOutsideAction.lean}
constructs the actual outside equivalence, proves
rotation-equivariance from genuine EPPA lifts,
and derives the four-form reduction and five-edge
conclusion. The principal final declaration is
\texttt{cycleFive\_eppa\_proper\_outside\_two\_forms}.
Certification requires the consolidated Lean CI
and axiom audit; this note does not claim the
full Lemma~3.18 is formalised.

\paragraph{Suggested one-paragraph substitution.}
After introducing outside neighbourhood profiles,
use the orbit-budget argument to obtain their canonical
cyclic labels whenever a proper profile occurs;
then derive \(H[B]\cong C_5\) from the cyclic
pair-orbit classification and equal half-edge counts.
The all-empty/full case should be addressed separately.
\end{document}
